\begin{theorem}[SmoothedChebyshevDirichlet]
\uses{SmoothedChebyshev,Smooth1}
\label{SmoothedChebyshevDirichlet}
\lean{SmoothedChebyshevDirichlet}
% at /home/runner/work/trough/trough/library/PrimeNumberTheoremAnd/MediumPNT.lean:262.0-379.28
\leanok
We have that
  $$\psi_{\epsilon}(X) = \sum_{n=1}^\infty \Lambda(n)\widetilde{1_{\epsilon}}(n/X).$$
\end{theorem}
\begin{proof}
\uses{Smooth1Nonneg,Smooth1Properties_above,Smooth1ContinuousAt,SmoothedChebyshevDirichlet_aux_integrable,LogDerivativeDirichlet,Smooth1LeOne,SmoothedChebyshevDirichlet_aux_tsum_integral}
\leanok
We have that
  $$\psi_{\epsilon}(X) = \frac{1}{2\pi i}\int_{(2)}\sum_{n=1}^\infty \frac{\Lambda(n)}{n^s}
  \mathcal{M}(\widetilde{1_{\epsilon}})(s)
  X^{s}ds.$$
  We have enough decay (thanks to quadratic decay of $\mathcal{M}(\widetilde{1_{\epsilon}})$) to
  justify the interchange of summation and integration. We then get
  $$\psi_{\epsilon}(X) =
  \sum_{n=1}^\infty \Lambda(n)\frac{1}{2\pi i}\int_{(2)}
  \mathcal{M}(\widetilde{1_{\epsilon}})(s)
  (n/X)^{-s}
  ds
  $$
  and apply the Mellin inversion formula.
\end{proof}
